\subsection*{数学中的真理性概念}

数学家们始终坚信他们所证明的是“真的”。显然，这种信念只能属于情感或形而上学的范畴，在数学领域内无法被证明合理，甚至无法赋予其非同义反复的意义。因此，数学中真理概念的历史属于哲学史而非数学史；但这一概念的演变对数学的发展产生了不可否认的影响，正因如此，我们不能对其避而不谈。

首先，一位精通哲学的数学家，就像一位拥有广博数学知识的哲学家一样罕见。数学家对哲学问题的看法，即使是涉及他们自身科学的问题，也往往是二手或三手的，且来源价值可疑。但恰恰因此，这些“普通”观点对数学史家的兴趣，不亚于笛卡尔或莱布尼茨（仅举两位也是最高级别数学家的思想家）、柏拉图（至少熟悉他那个时代的数学）、以及亚里士多德或康德（对后者则不能这样说）等思想家的原创见解。

传统的数学真理观念可以追溯到文艺复兴时期。在这种观念中，数学所研究的对象与自然科学所研究的对象之间没有重大区别。两者都是可知的，并且可以通过直觉和理性来把握；只要恰当运用，直觉和理性都不会出错。帕斯卡说：“只有完全错误的理智，才会从如此明显以至于几乎不可能被忽视的原理中错误地推理。”{[}11{]}，第XII卷，第9页）。笛卡尔确信“只有数学家才能找到证明，也就是说，确定而明显的理由”（{[}10{]}，第VI卷，第19页），而这一点（如果我们接受他的说法）是在他构建形而上学之前就已确立的，他在该形而上学中说：“这一点，我刚才将其作为规则，即我们非常清楚且非常分明地构想的事物都是真的，仅仅因为上帝存在或存在且是一个完美的存在者而得到保证”（{[}10{]}，第VI卷，第38页）。尽管莱布尼茨反对说，如何识别一个观念是“清楚且分明的”（*）并不明显，但他也认为公理一旦术语被理解，就是定义的自明且不可避免的推论（†）。此外，不应忘记，在这一时期，数学涵盖了许多科学——有时甚至包括我们今天不视为数学的工程师技艺。其在“自然哲学”、“机械技艺”和航海中的应用所取得的显著成功，是它激发信心的主要原因。

从这个观点来看，公理与逻辑推理的规则一样，不容讨论或怀疑；至多，一个人是“以古人的方式”进行推理，还是放任自己的直觉，这属于个人选择的问题。出发点的选择同样是个人的偏好。因此，欧几里得著作的许多版本得以出版，在这些版本中，《几何原本》坚实的逻辑框架被奇怪地歪曲了，并且给出了微积分或理性力学的“演绎”论述，而其基础却铺设得极其糟糕。因此，斯宾诺莎声称他的《伦理学》是“以几何学家的方式证明的”（more geometrico demonstrata）时，或许确实是出于真诚。尽管在十七世纪很难找到两位数学家在任何单一问题上意见一致，尽管争论频繁、无休无止且激烈尖刻，然而真理的概念却免于一切责难。“既然每一事物只有一种真理”，笛卡尔说，“无论谁找到它，就知道了一切可知之事”（{[}10{]}, 第VI卷，第21页）。

尽管关于这些问题的伟大时期的希腊数学文本没有保存下来，但希腊数学家在这个问题上的观点很可能远没有那么僵化。推理规则只是通过经验才发展到足以激发完全信心的地步；在它们能被视作超越一切讨论之前，必定经历了无数的犹豫和尝试性的形式。此外，如果帕斯卡认为不言自明（并且根据他姐姐所创的传说，他自己在孩童时期就以万无一失的直觉发现了这些公理）的那些“公理”没有成为长期讨论的对象，那将与希腊人的批判精神以及他们对讨论和诡辩的爱好格格不入。在一个严格来说并非几何学的领域中，埃利亚学派哲学家的悖论为我们保留了一些此类争论的痕迹；而阿基米德在观察到（{[}4{]}, 第265页）他的前辈们曾多次使用以他名字命名的公理时，补充说，借助这条公理所证明的东西“与不借助它证明的东西相比，其被接受的程度毫不逊色”，并且对他来说，他自己的结果能在此基础上被接受就足够了。柏拉图依照他的形而上学观点，将数学呈现为通往“自在真理”的途径，并将其所处理的对象视为在理念世界中独立存在。他在《理想国》的一段著名文字中准确地刻画了数学方法：“那些研究几何和算术的人\ldots{} 假定奇数和偶数以及三种角的存在；这些事物他们视为已知，并认为无论对自己还是对他人，都无需加以证明，因为它们对每个人来说都是不言自明的；并且，从这些出发，他们前后一致地一步步推进到他们着手考察的命题”（第六卷，510c-e）。因此，构成证明的首先是某种出发点，它在某种程度上是任意的（尽管“对每个人来说都是不言自明的”），并且，他稍后说道，人们不会试图再追溯得更远；其次，是一条依次经过一系列中间站的路径；最后，在每一步中，对话者的同意保证了推理的正确性。还应当补充的是，一旦公理被陈述出来，原则上就不再允许新的直觉诉求；普罗克洛斯引用格米努斯的话回忆道：“我们甚至教导这门科学的先驱们，在涉及将成为我们几何学说一部分的推理时，不要考虑仅仅貌似可信的结论”（{[}3{]}, 第I卷，第203页）。

因此，数学推理的规则是在经验与批评的烈火中发展起来的。如果正如有人颇有道理地论证的那样（*），欧几里得第八卷为我们保存了阿尔基塔斯算术的一部分，那么在那里发现那种学究式的刻板推理也就不足为奇了，这种刻板在任何发现或自以为发现了“严谨”的数学学派中都绝不会缺席。然而，一旦这些推理规则进入了数学家的共同实践，它们似乎直到很近的时代才受到质疑；尽管亚里士多德和斯多葛学派通过推理图式从某些规则推导出另一些规则，原始规则始终被视为不证自明。同样，在回溯到那些在他们看来为其时代科学提供了坚实基础（例如，大约公元前450年，传统上归于希俄斯的希波克拉底的《原本》初版中必定出现的那样）的“假设”、“公理”和“公设”之后，古典时期的希腊数学家似乎将精力倾注于发现新结果，而非对基础进行批判，因为在这一时期，这种批判只可能是徒劳的；并且，撇开形而上学的关切不谈，前面引用的柏拉图文本证明了数学家们对其科学基础存在这种普遍共识。此外，希腊数学家似乎并不认为作为他们出发点的“基本概念”——直线、面、量的比——能够进一步阐明。如果他们给出了“定义”，那显然是为了心安，而对这些定义的有效性不抱任何幻想。不言而喻，相比之下，希腊数学家和哲学家对于除“基本概念”之外的定义（通常称为“名义”定义，起着与我们“缩写符号”相同的作用）有着完全清晰的认识。在这方面，数学中“存在性”的问题明确地介入，无疑这是第一次。亚里士多德没有忽略指出，定义并不蕴含被定义事物的存在，因此需要公设或存在性证明。毫无疑问，他的观察源于数学家的实践。在任何情况下，欧几里得都谨慎地公设圆的存在，证明等边三角形、平行线、正方形等的存在，正如他在论证中需要引入它们时所做的那样。{[}3{]}第一卷）；这些证明是“构造”——换言之，他借助其公理展示出数学对象，随后证明这些对象满足所涉定义。

因此我们看到，古典时期的希腊数学达到了一种经验上的确定性（无论其形而上学基础可能是什么，依据这位或那位哲学家的观点）。如果数学推理的规则被认为无可置疑，那么希腊科学的成功，以及认为批判性修正不合时宜的感觉，在很大程度上归因于公理所激发的信心：这种信心与上个世纪人们对理论物理学原理所给予的信任属于同一层次。这正是谚语“nihil est in intellectu quod non prius fuerit in sensu”所暗示的，而笛卡尔正确地对此提出异议，认为它未能为他打算通过运用理性所实现的目标提供足够坚实的基础。

直到十九世纪初，数学家们才从笛卡尔的傲慢立场（更不用说康德或黑格尔了；后者在与其时代科学的关系上，如所预料的那样，稍显落后（*））退回到像希腊人那样灵活的立场。对古典观念的第一击，是本世纪初高斯、罗巴切夫斯基和波利亚伊构造了双曲非欧几何。我们在此不打算详述这一发现的起源，它是无数次试图证明平行公设的徒劳尝试的顶峰。当时，它对数学原理的影响或许并不像有时所断言的那样深刻。它只是要求放弃前一个世纪对欧几里得几何“绝对真理”的主张，并且理所当然地放弃莱布尼茨关于定义蕴含公理的观点；后者不再显得是“自明的”，而更像是假设，要根据它们是否适合物理世界的数学表示来判断。高斯和罗巴切夫斯基相信，各种可能几何之间的争论可以通过实验来解决（{[}15{]}, 第76页）。这也是黎曼的观点，他那篇著名的就职演讲《论几何学基础所依据的假设》的目的，就是为各种自然现象提供一个普遍的数学框架。“问题仍有待解决，”他说，“这些假设在多大程度上以及以何种方式得到实验的证实”（{[}19{]}，第284页）。但显然，这是一个与数学不再有任何关系的问题；而且刚才提到的那些作者中，似乎没有人怀疑过，即使一种“几何”不符合实验现实，其定理在“数学上”也不因此就不那么“为真”（*）。

然而，倘若果真如此，这种信念便不能再归因于对经典“几何直观”的无限信赖。黎曼试图为其著作对象——“n 重延展的流形”——所给出的描述，仅仅在证明引入“局部坐标”的合理性这一点上依赖于“直观”的考量（†）；自此以后，他似乎觉得自己已立足于坚实的地基，即分析学的地基。但分析学最终建立在实数概念之上，而实数概念直到那时仍停留在十分直观的基础上；函数论的发展正导致在这方面出现最令人不安的结果。随着黎曼本人关于积分的工作，尤其是博尔扎诺和魏尔斯特拉斯所构造的无切线曲线的例子，数学的全部病态现象开始显现。一个世纪之后，我们已经见过如此之多的这类怪物，以至于对此已有些麻木。但在十九世纪大多数数学家身上所产生的效果，则从厌恶到惊骇不一而足。“直观何以在此处欺骗我们？”庞加莱问道（{[}33 b{]}，第19 页）；而埃尔米特（不无幽默地，这一点并非所有评论这句名言的人都似乎察觉到了）宣称，他“怀着恐惧与惊骇，避开了这种没有导数的连续函数这一可悲的祸害”（{[}27{]}，第II 卷，第318 页）。这种不安最严重的特征在于：这些与常识如此相悖的现象，已不能再像“不可分量”时代那样归咎于尚未阐明清楚的概念，因为它们是在博尔扎诺、阿贝尔和柯西的改革之后出现的，而正是这些改革成功地将极限概念建立得与比例论一样严格。因此，它们只能归因于我们几何直观的粗糙和不完备，而自那时起，几何直观作为一种证明方法便已声名扫地。

这一结论不可避免地影响了经典数学，首先是几何学。无论人们对欧几里得的公理化构造怀有多大的敬意，其缺陷从古代起就已被注意到。平行公设曾是批评和证明尝试最多的对象；但欧几里得的后继者和评注者们也曾试图证明其他公设（特别是直角相等公设），或者认识到某些定义（如直线或平面的定义）的不充分。在 16 世纪，《原本》的编者克拉维乌斯注意到缺少一个保证第四比例项存在的公设；莱布尼茨则指出，欧几里得使用了几何直观却未明确提及，例如当他假定（《原本》第一卷命题 1）两个圆，每个都经过另一个的圆心，有一个公共点时（{[}12 b{]}，第VII 卷，第166 页）。高斯（他本人并不拒绝使用这类拓扑学的考量）提请人们注意，在欧几里得的作图中，位于另外两点（或两线）“之间”的点（或线）这一概念所起的作用，而这一概念却从未被定义（{[}14{]}，第VIII 卷，第222 页）。最后，位移的使用——尤其是在“全等三角形”的情形中——长期以来被认为理所当然（*），不久在十九世纪的批评者看来，乃是一个建立在未明确表述的公设之上的概念。因此，在 1860 年至 1885 年期间，出现了许多对几何基础的部分修订（亥姆霍兹、梅雷、乌埃尔），目的是填补其中的某些空白。但直到帕施 {[}28{]} 才明确地将放弃一切对直观的诉诸作为一个纲领提出，并以完全的严格性加以实施。他的事业的成功很快产生了众多效仿者，他们主要是在 1890 年至 1910 年间，给出了欧几里得几何公理的各种不同表述。这些著作中最著名的是皮亚诺的著作，用他的符号语言写成 {[}29 b{]}，尤其是希尔伯特的《几何基础》，该书于 1899 年问世，并以其阐述的清晰性和深刻性立即且当之无愧地成为现代公理化的典范，以至于其前辈们都被遗忘了。希尔伯特不满足于给出欧几里得几何的一个完整公理系统，而是按照其性质将它们分为若干组，并着手确定每组公理的精确范围，不仅通过孤立地发展每组公理的逻辑推论，而且通过讨论通过删去或修改某些公理而得到的各种“几何”（其中罗巴切夫斯基和黎曼的几何作为特例出现（*））。这样，他在一个直到那时仍被视为最接近外部实在的领域中，清晰地揭示了数学家在选择其公设时所拥有的自由。尽管这些具有怪异性质的“元几何”在部分哲学家中间引起了混乱，《基础》的论点却迅速且几乎一致地被数学家们所接受。因此，很难被怀疑偏向形式主义的庞加莱，于 1902 年断言，几何公理是约定，对于这些约定，“真”这一概念在日常意义上毫无意义（{[}33 a{]}，第66-67 页）。“数学真理”因此完全在于从任意设定为公理的前提进行逻辑演绎。正如我们稍后将看到的，演绎所依据的推理规则的有效性不久也遭到质疑，从而导致对数学基础之观念的一次彻底革新。

